% PIBA Template -- Software Architecture Lab, WS 2026/27, University of Stuttgart
%
% Self-contained: compile with pdflatex (no additional files needed).
% Rules: the whole PIBA must fit on ONE page. Do not change font size or margins.
% Replace all text in \todo{...} and remove the hints before submitting.

\documentclass[11pt,a4paper]{article}

\usepackage[utf8]{inputenc}
\usepackage[T1]{fontenc}
\usepackage{lmodern}
\usepackage[a4paper,margin=2cm]{geometry}
\usepackage{enumitem}
\usepackage{graphicx}
\usepackage{xcolor}
\usepackage[hidelinks]{hyperref}

\pagestyle{empty}
\setlength{\parindent}{0pt}
\setlength{\parskip}{0.5em}
\setlist{itemsep=0.1em, topsep=0.2em}

% Placeholder marker; leave the definition, but remove all uses before submitting.
\newcommand{\todo}[1]{\textcolor{gray}{[#1]}}

% Part heading
\newcommand{\pibapart}[1]{\vspace{0.4em}{\large\textbf{#1}}\par}

% ---------------------------------------------------------------------------
% Project information
% ---------------------------------------------------------------------------
\newcommand{\projecttitle}{PolyVote: A Declarative, Verifiable E-Voting Architecture}
\newcommand{\teamname}{Tokyo}
\newcommand{\teammembers}{David Enoghama, Franka Heger, Clemens Korneli, Patrick Ng, Felix Roehr, Jerome Stephan}

\begin{document}

\begin{center}
  {\LARGE\textbf{\projecttitle}}\\[0.4em]
  {\large PIBA -- Software Architecture Lab, WS 2026/27}\\[0.3em]
  \teamname\ \textbullet\ \teammembers
\end{center}

% ---------------------------------------------------------------------------
\pibapart{Problem}
Organizations running digital elections (university committees, student unions, associations) face a trade-off between \textbf{ballot flexibility} and \textbf{cryptographic verifiability}. Practical e-voting systems (\emph{Helios}, \emph{Belenios}) enforce rigid formats supporting only 1-of-$N$ or approval voting, while flexible research frameworks (\emph{Kryvos}, \emph{Ordinos}) lack usable implementations. Survey tools like \emph{LimeSurvey} let officials build custom GUI ballots, but lack verifiability and require blind trust in a central admin.

% ---------------------------------------------------------------------------
\pibapart{Idea}
We propose \textbf{PolyVote}, an e-voting architecture combining a \textbf{No-Code Election Studio GUI} and a \textbf{Declarative Election-as-Code Compiler Pipeline} around an encapsulated General-Purpose Zero-Knowledge Proof (GPZKP) adapter. Officials visually configure custom voting schemes (e.g., Ranked-Choice, Quadratic, or Cumulative Voting) in the Studio GUI, emitting a declarative JSON/DSL specification. Our Compiler Pipeline translates this specification into three synchronized artifacts: (1)~a dynamic \emph{Voter Ballot UI} schema, (2)~constraints for an abstract \emph{GPZKP Prover/Verifier Adapter}, and (3)~accumulation rules for the \emph{Voting Authority Service}. Voters compute encrypted ballots and GPZKP validity proofs client-side and post them to an append-only \emph{Bulletin Board} for verification.\\
\textbf{Quality Attributes \& Scope:} Driven by \emph{Modifiability \& Configurability} (new voting modes without code changes), \emph{Decoupling} (isolating the black-box GPZKP primitive behind an Adapter), and \emph{Auditability}. \emph{Out of Scope:} Cryptographic performance tuning, new ZKP primitives, Benaloh challenges.

% ---------------------------------------------------------------------------
\pibapart{Benefits}
\begin{itemize}
  \item \textbf{No-Code Ballot Expressiveness} --- addresses the rigid formats of \emph{Helios}/\emph{Belenios} and missing implementations of \emph{Kryvos}/\emph{Ordinos}; demonstrated by configuring and compiling Plurality, Ranked-Choice, and Quadratic Voting in the GUI without code changes.
  \item \textbf{Survey Usability with Public Auditability} --- addresses \emph{LimeSurvey}'s black-box trust model; demonstrated by posting encrypted ballots and validity proofs to the Bulletin Board.
  \item \textbf{Cryptographic Agnosticism \& Isolation} --- addresses tight coupling to experimental cryptography; demonstrated by swapping the GPZKP adapter between a mock and the real primitive.
\end{itemize}

% ---------------------------------------------------------------------------
\pibapart{Action}
\begin{enumerate}
  \item \textbf{Requirements \& Domain Models (27.11.\ \& 11.12.2026):} Elicit requirements (27.11.) and model the domain and declarative \emph{Election-as-Code} grammar for ballots and budgets (11.12.).
  \item \textbf{Architecture Design (08.01.2027):} Specify component boundaries, the three-target Compiler Pipeline, Bulletin Board protocol, and Voting Authority verification workflow.
  \item \textbf{Prototype \& Evaluation (22.01.--05.02.2027):} Build the \emph{Minimum Deliverable} (Studio GUI, Compiler, Voter Client, Bulletin Board, Voting Authority) and evaluate modifiability.
\end{enumerate}
\textbf{Tech \& Team Split (6):} Java 21/Spring Boot, TypeScript, PostgreSQL, Docker: (1)~\emph{Studio GUI}, (2)~\emph{DSL/Validator}, (3)~\emph{Compiler}, (4)~\emph{Voter/Prover Adapter}, (5)~\emph{Bulletin Board}, (6)~\emph{Voting Authority}.

% ---------------------------------------------------------------------------
\vfill
{\small\textbf{Prior work disclaimer.} One member worked on GPZKP ballot-validity primitives in their Bachelor's thesis; to avoid conflict of interest, it is used strictly as a black-box dependency without runtime guarantees.\\
\textbf{AI usage statement.} Google Gemini was used to brainstorm the architecture and condense the PIBA text.}

\end{document}